% Section 7 -- Limitations (Yang 2027 pattern: choice -> when it holds -> when it fails -> future work).
\section{Limitations}\label{sec:limitations}

The main limitation is that the class-pure sites were constructed from cohorts in which the sites hold both classes. Drawing both roles at random from the same cohort controls part, but not all, of the difference between case sites and control sites, and does not guarantee (A1$'$): a site is eligible for a role only if it has enough participants of that class, so some sites can serve only as case sites or only as control sites, and the largest sites never serve as test sites. In real collections, cases from a clinic and controls from a school may also differ in age, medication or motion, a bias that none of our methods removes. The natural collections show that the evaluation bias is real, but not whether bridging removes it there.

The theory assumes additive site effects (A0) and linear scores. Bridging and within-site training remove a site shift common to both classes but cannot address a site-by-diagnosis interaction, such as a scanner that distorts patients' signals differently. Proposition~\ref{prop:leak} was formulated after the experiments, and our L3 implementation is shift-invariant only approximately.

The constructed designs are also small, with $20$ to $40$ subjects per class and three to five test sites per repeat, so single repeats vary widely and LSPO remained imprecise (mean absolute error $0.09$ to $0.19$ across the two ABIDE cohorts). In ADHD-200, only $k=1$ could be constructed and the diagnostic signal was close to chance, so several criteria could be assessed only in part.

The comparison is also limited in scope: the non-transductive setting of the published methods required an extension of the authors' code, and ComBat was fitted once on the outer training set, so our harmonization result describes this pipeline only. ABIDE-II shares some institutions with ABIDE-I, although not participants, preprocessing or atlas. Finally, all features were resting-state fMRI connectivity; we have not tested the framework on EEG or structural MRI.
